\documentclass[11pt]{article}
\usepackage[utf8]{inputenc}
\usepackage{amsmath, amssymb, amsthm}
\usepackage[margin=1in]{geometry}
\usepackage{booktabs}
\usepackage{hyperref}

\title{Review of \texttt{scratch.tex}}
\author{}
\date{}

\begin{document}
\maketitle

This is a review of the in-progress derivation document
\texttt{scratch.tex} (and \texttt{main.tex}, the more polished
short-form version). Comments are organized by substantive technical
issues, smaller items, and writing.

\section*{Overall}

The project is sensible and worth doing: take Santiago--Caballero
(1995, 1998) plus Buffalo \& Kern (2024) plus a mating-system layer,
parameterize with empirical $V_A$ from wild-pedigree animal models,
and see how much $N_e$ depression is implied vs.\ what is observed.
The math is mostly tractable. The current draft has, however, (a) a
``master equation'' that is presented as exact but is actually a
leading-order approximation, with a missing $1/(1+\Lambda_u)$ in the
exponent that biases the headline numbers, (b) at least one derivation
step that is wrong as written ($\rho_{mf} = 1-p$), (c) a
unit/conventions problem with $V_A$ that runs through the whole
document, and (d) a fair amount of rhetorical hyperbole that should be
deleted on a global pass.

\section{Substantive technical issues}

\subsection{The ``Master Equation'' is an approximation, presented as exact, with a missing $1/(1+\Lambda_u)$ in the exponent}

Section 4 (lines 222--227):
\[
\frac{N_e}{N} = \frac{\exp(-\Lambda_{\text{linked}})}{1 + \Lambda_{\text{unlinked}}}
\]

The form is more defensible than my first read suggested. Starting
from a unified discrete expression $1/(1+\Lambda_l + \Lambda_u)$ and
pulling the linked piece out:
\[
\frac{1}{1+\Lambda_l + \Lambda_u}
= \frac{1}{(1+\Lambda_u)\big(1 + \tfrac{\Lambda_l}{1+\Lambda_u}\big)}
\;\approx\; \frac{1}{1+\Lambda_u}\,\exp\!\left(-\tfrac{\Lambda_l}{1+\Lambda_u}\right)
\;\approx\; \frac{e^{-\Lambda_l}}{1+\Lambda_u}.
\]
The first approximation uses $1/(1+u) \approx e^{-u}$ for small
$u = \Lambda_l/(1+\Lambda_u)$. The second drops the $1/(1+\Lambda_u)$
in the exponent, which requires $\Lambda_u$ small. So the mixed form
is the leading-order approximation to a unified
$1/(1+\Lambda_l+\Lambda_u)$ in the regime where $\Lambda_l$ is small;
the prose on lines 220--221 even invokes that condition (``the local
variance on a single haplotype is a small fraction of the genome-wide
total''). Two issues survive:

\paragraph{It's an approximation, not ``exact.''} Line 224 calls this
the ``exact Master Equation.'' The honest framing is: it's the
leading-order approximation to $1/(1+\Lambda_l+\Lambda_u)$ when
$\Lambda_l \ll 1+\Lambda_u$. The justification offered ---
``crossing-over and independent assortment are distinct molecular
mechanisms, their probabilistic effects on the survival of a neutral
lineage are strictly independent'' --- is overselling it; independence
of mechanisms doesn't make the multiplicative form exact, it makes it
a leading-order combination.

\paragraph{The cleaner second-order form is $\exp(-\Lambda_l/(1+\Lambda_u))/(1+\Lambda_u)$, not $\exp(-\Lambda_l)/(1+\Lambda_u)$.}
The factor of $1/(1+\Lambda_u)$ inside the exponent matters in exactly
the regime the stress test cares about. Numerically, for
$\Lambda_l = 0.1$ and $\Lambda_u = 2$:

\begin{center}
\begin{tabular}{lc}
\toprule
Form & $N_e/N$\\
\midrule
Unified discrete $1/(1+\Lambda_l+\Lambda_u)$ & $0.3226$\\
Cleaner approximation $\exp(-\Lambda_l/(1+\Lambda_u))/(1+\Lambda_u)$ & $0.3223$\\
Draft's $\exp(-\Lambda_l)/(1+\Lambda_u)$ & $0.3016$\\
\bottomrule
\end{tabular}
\end{center}

The draft's form is biased downward by $\sim 7\%$ in this regime ---
same order of magnitude as the $\kappa$ error from issue 2. Either
re-derive with the $1/(1+\Lambda_u)$ correction inside the exponent,
or work in the unified $1/(1+\Lambda_l+\Lambda_u)$ form, or use the
unified $\exp(-\Lambda_l-\Lambda_u)$ form throughout (as
\texttt{main.tex} does). All three are defensible; the current mixed
form sits between them and inherits the worst of each.

\subsection{$\rho_{mf} = 1-p$ is asserted, not derived, and I think it's wrong as stated}

Line 130: ``The statistical correlation between their reproductive
outputs therefore scales linearly with this within-pair paternity
share ($\rho_{mf} = 1 - p$).''

Set up the moments. Female $i$ has $K_i$ surviving offspring. A
fraction $1-p$ are sired by her social mate; $p$ are sired
extra-pair. Social male's reproductive success is
$k_m = (1-p)K + S$, where $S$ is his extra-pair siring success across
all females. Then
\[
\rho_{mf}
= \frac{(1-p)\,\mathrm{Var}(K)}{\sqrt{\big[(1-p)^2\mathrm{Var}(K) + \mathrm{Var}(S)\big]\,\mathrm{Var}(K)}}.
\]
This collapses to $1-p$ only under restrictive assumptions about
$\mathrm{Var}(S)$. If $S$ is Poisson with rate $p\bar K$ (mass-action
mating), $\rho_{mf}$ exceeds $1-p$. If extra-pair success is itself
heritable and correlated with within-pair success, it differs again.
The bottom line: $\rho_{mf}$ depends on the variance of extra-pair
siring success, not just on $p$. This is the keystone of the entire
$\kappa = 2-p$ result, so it cannot be a one-line assertion. We need
an explicit model of EPP siring distribution and a proper derivation,
with bounds for the worst- and best-case variance assumptions.

(Nunney 1993 and Caballero 1995 derive related but distinct forms;
cite them and show how this expression connects.)

A companion document working through this carefully is in
\texttt{rho\_mf\_derivation.tex} (in the same directory).

\subsection{The unlinked baseline under ``random mating'' doesn't match a direct individual-based derivation}

In \S 1.3.2, the pair-based derivation gives $N_e = N/(1+\kappa V_A)$
with $\kappa=1$ recovering ``random mating,'' i.e.\ $N_e = N/(1+V_A)$.

But: under genuinely non-pair-bonded random mating, there are no
persistent pairs. An individual's gamete contribution has variance
$V_{km} = 2 + 4V_A$ (Poisson plus heritable, on the absolute-offspring
scale). Plugging into Eq.~8 ($N_e = 4N/(V_{km}+V_{kf})$) gives
\[
N_e = \frac{4N}{4 + 8V_A} = \frac{N}{1+2V_A},
\]
not $N/(1+V_A)$. There is a factor-of-two discrepancy. The pair
framing is silently changing the model when $\kappa=1$. Either the
pair framework is the wrong abstraction for non-monogamous systems,
or the ``random-mating'' limit needs a different scalar. Whichever it
is, the current derivation does not recover the standard
sex-separated Wright--Fisher random-mating result, and that is a
problem because the EPP scaling claims to interpolate between those
endpoints.

\subsection{Unit confusion on $V_A$ runs throughout}

$V_A$ is used interchangeably for:
\begin{itemize}
\item variance in \emph{relative} fitness $w = k/\bar k$ (Sections 1, 2, the master equation)
\item variance in \emph{absolute} fitness in the mutation--selection--drift derivation ($V_A \approx Us$, $\Delta V_{\text{mut}} = Us^2$)
\end{itemize}
These differ by $\bar w^2$ (or, on the absolute scale, by
$\bar k^2 = 4$). For populations near demographic equilibrium with
$\bar w \approx 1$ this is harmless; for load-heavy populations where
$\bar w = e^{-U}$ is materially less than 1 it isn't. The whole
``stress test'' rests on plugging empirical pedigree-derived $V_A$
(typically reported as variance in \emph{relative} fitness, sometimes
annual reproductive success, sometimes lifetime --- units vary by
study) into formulas that mix conventions. A dedicated paragraph at
the start of \S 4 nailing down conventions, and a check that the
empirical estimates are being plugged in consistently, is essential.

\subsection{The $v_i/2$ ``homologous partition'' deserves more care}

\S 3 claim: a focal neutral allele on chromosome $i$ experiences
linked selection only with the $v_i/2$ on its own physical molecule,
and the homologous $v_i/2$ enters as an independently-assorting
unlinked contribution.

The intuition is right but the formalism is hand-wavy. The $v_i$ as
derived in \S 3.1 is the population-level additive variance (summed
across both haplotypes and all individuals' diploid genotypes).
Splitting that into ``cis to focal'' and ``trans to focal'' needs an
explicit step: at any selected locus on chromosome $i$, the focal
neutral is on the same physical molecule as a \emph{given} deleterious
allele copy with probability $1/2$ --- and even then, ``the variance
linked in cis'' isn't simply half the population-level additive
variance of that locus, because the cis/trans distinction interacts
with how the variance was defined (per-allele effect vs.\ per-genotype).
I'd want to see this written out explicitly, or at minimum a citation
to where Buffalo \& Kern (2024) formalize it. As is, it's a gesture.

\subsection{The $V_A$ derivation in \S 4 conflates two conventions}

The draft writes $V_m = Us^2$ and $\Delta V_{\text{drift}} = -V_A/(2N_e)$
and $\Delta V_{\text{sel}} = -sV_A$, then solves
$\hat V_A = 2N_eUs^2/(1+2N_e s)$. This is the standard
infinite-sites-additive result and matches Hill (1982) and
Charlesworth (2015) in the appropriate convention. But:
\begin{itemize}
\item Whether $U$ is haploid or diploid changes the answer by 2.
\item Whether $s$ is the heterozygous selection coefficient or the
homozygous one matters for $\Delta V_{\text{sel}} = -sV_A$ (this
assumes purely additive, $h=1/2$ in $1, 1-s, 1-2s$ notation, with
$s$ heterozygous).
\end{itemize}
Both conventions need to be stated up front. The \S 4.4 ``Validation
via Diffusion'' paragraph (line 346) claims the diffusion integration
``converges \emph{exactly}'' to the heuristic. It does not, in general
--- this is a leading-order match, and the stationary density it
writes ($f(q) \propto e^{-4Nsq} q^{4N\mu - 1}(1-q)^{-1}$) has a
$(1-q)^{-1}$ factor that is improper at $q=1$, suggesting either a
typo or a careless transcription of the irreversible-mutation Wright
distribution. Needs to be fixed.

\subsection{The \S 2.4 high-recombination weak-selection limit doesn't match \texttt{main.tex}}

\texttt{scratch.tex} (lines 282--286) gives
$E_{\text{linked}} \approx \kappa v_i^2/(MU_i s^2)$; in the
weak-selection regime $v_i \approx 2 N_e U_i s^2$, so
$E_{\text{linked}} \approx 4\kappa N_e^2 U_i s^2 / M$.

\texttt{main.tex} (line 119) gives
$E_{\text{linked(weak)}} \approx 2\kappa U_i N_e \overline{s_h}/(3M)$
--- i.e., a different equilibrium-variance prefactor
($\tfrac{2}{3} U_i N_e s^2$, presumably from a Hudson--Kaplan-style
derivation).

These two are inconsistent. Two different $V_A$ equilibrium
expressions are floating around between the drafts. Pick the one
whose derivation can be defended, and discard the other.

\subsection{Self-consistency: $N_e$ enters on both sides}

$Z = 1 - V_m/V_A$ in the chromosomal integral, and $V_A$ at MSDE
depends on $N_e$, but $N_e$ is what is being predicted. The framework
is implicitly self-consistent: predicted $N_e$ and equilibrium $V_A$
should solve simultaneously. The derivation should either solve it
explicitly or note that \emph{empirical} $V_A$ is being plugged in
(which is what the draft says it is doing --- but then the ``stress
test'' is asking whether the empirical $V_A$ is consistent with
\emph{any} $(U, s, N_e)$ triple under MSDE, and that is a different
framing than the document currently makes clear).

\subsection{The ``paradox'' needs to engage harder with environmental confounding}

The conclusion (``massive empirical $V_A$ cannot be standard
equilibrium load, necessitating \dots\ balancing selection,
fluctuating optima, or systemic environmental confounding'') is the
right shortlist, but in the wild-pedigree literature the dominant
explanation for inflated $V_A(w)$ is well-known to be unmodeled
environmental and parental effects (cohort, territory quality,
maternal effects bleeding into the additive component because of
incomplete pedigrees). Bonnet et al.\ 2022 itself and the surrounding
literature (Kruuk, Postma, Wilson) discuss this. The paper as
currently framed treats environmental confounding as one possibility
among several at the very end; if that is actually the most likely
answer, the framing of ``theoretical crisis'' overstates the puzzle.
Engage with the animal-model literature explicitly and stake out a
position.

\section{Smaller things}

\begin{itemize}
\item Line 34: stray \verb|."| after the citation --- leftover quotation mark.
\item Line 338: the appendix subsection (``Derivation of Standing
Variance at MSDE Balance'') is one massive run-on with backslashes and
\verb|\u| artifacts where line breaks got lost --- looks like a paste
from another tool. Needs to be reformatted into proper LaTeX blocks.
\item Line 346: ``(cite your colleague / Hernandez et al.\ 2026)''
--- placeholder still in the text.
\item Section 5 (line 322) refers to ``the Delta method expansion
derived above,'' but the Delta-method derivation is in
\texttt{main.tex}, not \texttt{scratch.tex}. The two drafts need to be
unified or the cross-reference fixed.
\item The Robertson (1961) framing in \S 1 says $V_A$ ``captures the
independent, transmissible effects of alleles---assuming dominance and
epistatic interactions are negligible.'' Strictly, $V_A$ is
\emph{defined} as the additive component regardless of whether
dominance/epistasis are negligible; what is negligible is the
transmissible \emph{non}-additive variance via parent-offspring
covariance. Small wording fix.
\item Line 51 simplification: $1/(0.5 + 0.5s) = 2/(1+s)$ then ``$s$
negligible relative to 1.'' This is fine, but earlier (line 27) the
draft says ``agnostic to whether the fitness variance is driven by
advantageous sweeps or deleterious load.'' That is not quite true
once you take $s\to 0$ --- the framework recovers the classical S\&C
only under weak selection at \emph{each} selected locus. For sweeps,
$s$ may not be small. Worth flagging as a scope limit.
\end{itemize}

\section{Style / writing}

The draft is studded with words doing rhetorical work that should not
be there: \emph{exactly}, \emph{exact}, \emph{strictly},
\emph{perfectly}, \emph{rigorously}, \emph{elegantly},
\emph{profoundly simple}, \emph{highly rigorous improvement},
\emph{flawlessly captures}, \emph{mathematically proves},
\emph{definitive test}. Most of these should be deleted on a global
pass. ``Math proves X'' should become ``Eq.~(n) shows X'' or ``under
assumptions Y, X follows.'' Right now the prose carries a confidence
the derivations don't yet earn.

The document also tends to dramatize transitions (``To formally
quantify\dots'', ``To rigorously establish\dots''). A scientific
reader will find this exhausting after the first few sections.
Compress.

\section{Suggested order of fixes}

\begin{enumerate}
\item Either derive $\rho_{mf}$ properly with an explicit EPP siring
model, or replace $\kappa = 2-p$ with bounds (issue 2). See
\texttt{rho\_mf\_derivation.tex} for one explicit treatment.
\item Reconcile the random-mating limit (issue 3). Probably means
rewriting \S 1.3 around individuals and adding the pair-correlation
as a perturbation rather than as the primary frame.
\item Stop calling the master equation ``exact''; either move to the
cleaner $\exp(-\Lambda_l/(1+\Lambda_u))/(1+\Lambda_u)$ form or pick
one approximation (all-\texttt{exp} or all-\texttt{1/(1+x)}) and use
it throughout (issue 1).
\item Add a ``Conventions'' paragraph nailing down what $V_A$, $U$,
$s$ mean each time they appear (issue 4).
\item Either drop \S 4.4's exactness claim or write the diffusion
calculation in full and check it (issue 6).
\item Reconcile \texttt{scratch.tex} and \texttt{main.tex}'s
weak-selection limits (issue 7).
\item Engage with the wild-pedigree environmental-confounding
literature in the discussion (issue 9).
\end{enumerate}

\end{document}
